\section{Experimental Setup}
\subsection{Study Data and Provenance}
The locked window is 19 February 2025--7 October 2026. The metadata-only
sample selected 181 stations from an eligible catalogue of 481 stations in
255 parsed cities. Full-window counts in Table~\ref{tab:data} are inherited
descriptive metadata; final-test labels are not opened to produce this draft.
All analysis discussed here is development-only.
\begin{table}[t]\centering\footnotesize
\caption{Measured dataset and current operational alignment. Completed measured targets are preserved; archive coverage is partial.}
\label{tab:data}
\setlength{\tabcolsep}{4pt}
\begin{tabular}{lr}
\toprule
Measure & Count\\\midrule
Usable measured stations / cities & 179 / 100\\
Retained measured station-days & 107,280\\
Valid / rejected targets & 64,799 / 42,481\\
Censored Severe targets & 383\\
Guarded CAMS-aligned development rows & 7,560\\
Aligned stations / cities / dates & 180 / 100 / 42\\
Aligned valid / rejected targets & 5,506 / 2,054\\
Verified raw archive chunks & 6 / 36\\
\bottomrule\end{tabular}
\end{table}
\subsection{Chronological and Spatial Splits}
Expanding-window validation begins with 90 training dates, then 30 calibration
dates and 30 evaluation dates, advancing by 30 dates. One-day embargoes
separate adjacent partitions. Every station's row for a date moves together.
Station eligibility requires 120 valid training targets, so the first 90-day
training fold is necessarily insufficient; later outcomes cannot rescue it.
The latest 90 dates, 10 July--7 October 2026, remain sealed.

Spatial validation trains through 31 October 2025 and tests from 2 November
2025 through the final pre-test day. Ten-fold 25-km urban-cluster and station
holdouts, and six region holdouts, distinguish unseen cities from unseen
stations in seen clusters. Local-history and monitor-free inference are
reported separately. The all-station reference uses the same dates and test
pairs. Unavailable held-out monitor feeds must be masked in the corresponding
inference regime.

\subsection{Baselines and Training-only Tuning}
Climatology uses the station's training targets in the same month, with
station and global means as fallbacks. Persistence uses the rolling measured
24-hour index ending at the cutoff. Raw CAMS uses the same AQI engine.
The primary exact comparison requires all three baselines on identical pairs;
secondary comparisons relax only persistence availability. Severe days remain
in their category and lower-bound diagnostics.

Hyperparameters and early stopping use the last 30 dates of a fold's training
window as inner validation. Calibration and evaluation labels do not tune the
learner. All seeds are 42. The carry-forward model has the lowest mean
rolling-fold MAE, with ties favouring Ridge and residual formulations; raw
CAMS is retained if no learned model improves that criterion. Ablation fixes
the learner while varying feature groups. Fire admission requires a favourable
95\% paired date-block interval.
\begin{table}[t]\centering\footnotesize
\caption{Declared primary search space. These are preregistered candidates, not selected parameters.}
\label{tab:search}
\setlength{\tabcolsep}{4pt}
\begin{tabular}{lp{.54\columnwidth}}
\toprule
Learner & Training-only choices\\\midrule
Ridge & $\alpha\in\{0.1,1,10,100\}$\\
LightGBM & Leaves 15/31; minimum child samples 20/60\\
XGBoost & Depth 4/6; minimum child weight 5/20\\
Boosting controls & Rate 0.05; at most 600 rounds; early stop 40; row and column fractions 0.8\\
\bottomrule\end{tabular}
\end{table}
\subsection{Metrics and Statistical Comparison}
Exact-target MAE is primary; RMSE, signed forecast bias, coefficient of
determination and six-category diagnostics are secondary. Poor-or-worse and
Very-Poor-or-worse detection, false alarms and lower-bound Severe shortfall
are reported explicitly. Regional, seasonal and city-size subgroups below
30 pairs retain counts but withhold inferential metrics. Intervals are assessed
at 80\% and 90\% nominal coverage, with width and severity-specific coverage.

Paired comparisons preserve station/date pairs. Moving seven-day date blocks
with 5,000 resamples provide percentile intervals; daily mean loss differences
use a Diebold--Mariano comparison with Newey--West variance and Holm adjustment
\cite{diebold1995comparing}. An interval spanning zero does not establish a
nonzero difference. The immutable non-CAMS study uses a separate 30-valid-day
station rule and fixed Ridge $\alpha=10$; it cannot select the primary model.
